%% ============================================================
%%  §10b  Agentic Development: AI-Assisted Formal Verification
%% ============================================================

\section{Agentic Development of the Verification Library}
\label{sec:agentic}

lean-pl-verify was developed with sustained assistance from Claude Code
(Anthropic's AI coding agent, powered by Claude Sonnet~4.6).
This section documents what the agent automated, what required human judgment,
and what this experience implies for AI-assisted formal verification.
We report this not as a claim that AI ``wrote'' the proofs, but as an honest
account of the human--AI division of labor in a non-trivial verification project.

\subsection{What the Agent Automated}

\paragraph{Boilerplate proof generation.}
The large majority of the 258 theorems follow one of four proof shapes
(§\ref{sec:rust-proofs}).  Once the first theorem of each shape was written
by hand and verified, the agent generated the remaining theorems in that
shape by pattern-matching on the function body and instantiating the template.
This was most valuable for:

\begin{itemize}[leftmargin=2em]
  \item The 48 \texttt{ElabSpec.lean} theorems: 16 functions × 2--4 properties each.
  \item The 47 \texttt{CharonSpec.lean} theorems: generated from the
        Charon-extracted defs using the \texttt{charon\_simp} tactic macro.
  \item The 33 \texttt{TypeScript/ElabSpec.lean} theorems: mirroring the
        Rust proofs for the TypeScript evaluator.
  \item The 6 \texttt{TypeScript/BugDetection.lean} theorems: bug-expose and
        fix-verify proofs generated by instantiating the same \rfl template.
\end{itemize}

\paragraph{Simp set construction.}
Identifying which simp lemmas are needed to reduce a given evaluator call is
tedious but mechanical.  The agent iterated on failing simp calls by reading
error messages, adding missing lemmas (\texttt{getElem?\_pos},
\texttt{List.getElem\_cons\_succ/zero}, \texttt{decide\_eq\_true/false}),
and verifying the fix.

\paragraph{The Charon pipeline (\texttt{charon2lean.py}).}
The 461-line Python translator from Charon JSON to Lean \texttt{LLBCFunDef}
terms was written almost entirely by the agent.
The agent handled the flat-list-to-CPS conversion, the AddChecked/SubChecked
desugaring, the overflow-check Assert elision, negative-integer literal
wrapping, and type deduplication — each discovered by reading the 104\,KB
LLBC JSON and adapting to Lean's parsing requirements.
Human involvement was limited to approving the overall approach and
unblocking specific Charon installation issues (the \texttt{rustc-dev}
component requirement).

\paragraph{Adequacy proof scaffolding.}
The 14-theorem adequacy proof in \texttt{Adequacy.lean} was largely
scaffolded by the agent after the human specified the top-level theorem
statement.  The agent proposed the induction structure, identified the
\texttt{split at h} and \texttt{change} patterns needed for do-notation
reductions, and handled the 10+ statement cases.

\subsection{What Required Human Judgment}

\paragraph{Architectural decisions.}
The core choices — shallow monadic embedding, fuel-based evaluation,
forward-backward pairs for \texttt{\&mut T}, the \texttt{ProgramSpec}
inductive, positional environments for TypeScript — were made by the human
author based on the research goals and prior PL verification experience.
These decisions cannot currently be automated because they require understanding
the paper's contribution claims and the mathematical elegance trade-offs.

\paragraph{Sorry elimination.}
The \texttt{prodRange\_factorial} and \texttt{factSafe\_implies\_ov}
sorry closures required human-directed proof strategy (right-peel lemma,
\texttt{interval\_cases} for factorial bounds).
The agent could fill in the tactic steps once the strategy was specified,
but could not discover the strategy independently.

\paragraph{TypeScript AST design.}
The decision to add \texttt{set\_} (mutable update) to \texttt{TSStmt}
rather than use de Bruijn rebinding tricks for loop variables required
understanding the proof implications of each choice.  The agent proposed
both options but deferred to the human on which produced cleaner proofs.

\paragraph{Paper structure.}
The high-level narrative, the choice of which theorems to highlight, and
the framing of the unification claim as a first-class contribution were
human decisions.

\subsection{Workflow and Persistent Memory Model}

\begin{figure}[h]
\centering
\small
\begin{tabular}{@{}ll@{}}
\toprule
\textbf{MEMORY.md entry type} & \textbf{Example} \\
\midrule
Proof pattern & \texttt{split at h} for nested match reductions \\
Known simp lemma & \texttt{List.getElem?\_cons\_zero} (not \texttt{getElem?\_pos}) \\
Architectural fact & FBPair encodes \texttt{\&mut T} via forward/backward pair \\
Build fix & \texttt{charon2lean.py}: op\_name can be dict \texttt{\{"Rem":"UB"\}} \\
Sorry status & \texttt{FactInvariant.lean}: 2 sorry → 0 sorry (resolved) \\
\bottomrule
\end{tabular}
\caption{Categories of entries in \texttt{MEMORY.md} (persistent cross-session state).}
\label{tab:memory}
\end{figure}

The development proceeded in a persistent conversation model: each session
began with the agent reading \texttt{MEMORY.md} (maintained via the
\texttt{Write}/\texttt{Edit} tools), which accumulated architectural facts,
proof patterns, discovered simp lemmas, and open task status.

\textbf{Error correction taxonomy.}
Table~\ref{tab:errors} categorises the types of errors the agent successfully
corrected versus where it needed human redirection.

\begin{table}[h]
\centering
\small
\caption{Agent error handling in lean-pl-verify development.}
\label{tab:errors}
\begin{tabular}{@{}lll@{}}
\toprule
\textbf{Error type} & \textbf{Outcome} & \textbf{Example} \\
\midrule
Missing simp lemma & Agent resolves & Added \texttt{getElem?\_cons\_succ} from error \\
Wrong fuel bound (off by 1) & Agent resolves & Increased fuel from 10 to 12 \\
Unicode in listings (UTF-8 byte) & Agent resolves & Added \texttt{lstdefinelanguage\{TypeScript\}} \\
Dict-form BinOp in Charon JSON & Agent resolves & \texttt{isinstance(op\_name, dict)} check \\
\midrule
Proof strategy (which lemma to use) & Human directs & Right-peel for \texttt{prodRange\_factorial} \\
Architectural choice & Human decides & \texttt{set\_} vs de~Bruijn rebinding for TS \\
Sorry closure strategy & Human directs & \texttt{interval\_cases} for factorial bounds \\
Context loss (re-proposing failed fix) & Human redirects & \texttt{MEMORY.md} partially mitigates \\
\bottomrule
\end{tabular}
\end{table}

Observed productivity gains:
\begin{itemize}[leftmargin=2em]
  \item \textbf{Build-debug cycles} were 3--5$\times$ faster: the agent
        read error messages, proposed fixes, and rebuilt without waiting for
        human interpretation of Lean's error output.
  \item \textbf{Repetitive proof patterns} (same proof shape × 16 functions)
        were handled in a single agent invocation rather than 16 manual edits.
  \item \textbf{Cross-file consistency} was maintained by the agent reading
        both the definition file and the proof file before suggesting changes.
  \item \textbf{Context-window management}: sessions compressed prior context
        automatically; \texttt{MEMORY.md} was the primary mechanism for
        preserving architectural decisions across compression boundaries.
        Structured proof-attempt logs (recording which approaches failed and why)
        would further improve multi-session continuity.
\end{itemize}

\subsection{Implications for Formal Verification}

The experiment suggests a productive division of labor:
humans specify the \emph{what} (theorem statements, proof strategies,
architectural choices) and agents handle the \emph{how} (tactic sequences,
simp sets, boilerplate generation, pipeline scripts).
This mirrors the relationship between a mathematician and a proof assistant
— the mathematician specifies the proof sketch; the assistant checks
the details — but with the agent absorbing an additional layer of mechanical
work above the kernel.

The main limitation observed was \textbf{context loss}: in long sessions,
the agent occasionally re-proposed solutions already known to fail,
requiring the human to redirect.
The persistent memory file (\texttt{MEMORY.md}) mitigated but did not
eliminate this.  Structured proof-attempt logs (recording which approaches
failed and why) would be a valuable engineering improvement.

We anticipate that agentic formal verification will become standard practice
as LLMs improve at understanding proof error messages and tracking proof
state across files.  lean-pl-verify demonstrates that the current generation
of agents can already handle a non-trivial 258-theorem library with
\textbf{0 sorry}.
